\documentclass[11pt]{article} 

\usepackage{amsmath}

\usepackage[super]{nth}

\usepackage{epsfig}

\usepackage{graphics}

\usepackage{wasysym}

\usepackage{times}

\usepackage{natbib}

%\usepackage{amssymb}

\setlength{\topmargin}{0in}
\setlength{\textwidth}{6.25in}

\setlength{\oddsidemargin}{.125in}

% 28 lectures total

\usepackage{amscd}

\begin{document}
\begin{center}

\large
\textsc{Mathematical Foundations of Political Methodology: PLSC 43401}\\
\end{center}

\small

\ \\ \begin{tabular} {ll} \multicolumn{2}{l}{Professor Robert Gulotty}\\
Office: &Pick 326\\
Phone: &2-5726\\
Email: &gulotty@uchicago.edu\\
\multicolumn{2}{l}{Office Hours: Tuesday/Thursday signup https://www.wejoinin.com/sheets/ppbii}
\end{tabular}

\vspace{.15in}
\vspace{.15in}

\ \\ \begin{tabular} {ll} \multicolumn{2}{l}{TA Arthur Yu}\\
Office: & Pick 507 \\
%Phone: & \\
Email: &stevenboyd@uchicago.edu  \\
%\multicolumn{2}{l}{R Lab: }\\
\multicolumn{2}{l}{Office Hours: Friday 4:30pm to 5:30pm}
\end{tabular}

\vspace{.5in}

\ \\{\bf Course description}  

\ \\This is a first course on the theory and practice of mathematical methods in social science research.  These mathematical and computer skills are needed for the quantitative and formal modeling courses offered in the political science department, and are increasingly necessary for courses in American Politics, Comparative Politics, and International Relations.   We will cover mathematical techniques (linear algebra, calculus, probability) and methods of logical and statistical inference (proofs and statistics).  

\vspace{.1in}

\noindent {\bf Grading:} You will be graded on 10 weekly problem sets (worth 40\% of your total grade), occasional quizzes covering the reading (10\%), and two, noncumulative exams (each worth 25\% of your grade).  

\ \\{\bf Problem sets are due each Tuesday at midnight, starting October $\mathbf{8^{\text{th}}}$.  Instructions for submission are as follows}

\begin{itemize}
\item Each student, ideally working with a group, works through the problem sets without looking at the answers and writes up their first attempt (can be handwritten)
\item Each student will, again as a group, discuss the correct answers and check.
\item Using the correct answers, students will correct their written responses from step 1, explaining their errors, what they were missing, or describing what they still do not understand about the solution (this is written in Latex).
\end{itemize}
Grades are evaluated on the basis of clarity more than how correct they managed to get.  The grader is your reviewer.


\ \\{\bf Required books:} \textit{ Mathematics for Economists, \nth{4} Edition},  by Pemberton \& Rau, and \textit{Introduction to Probability, \nth{2} Edition}, by Bertsekas \& Tsitsiklis.


%\newpage

\ \\ \noindent \textsc{Oct 1: Algebra, Sets, and Functions}

\begin{itemize}
\item Pemberton \& Rau Ch, 1-4

\end{itemize}




\noindent \textsc{Oct 3: Mathematical and Social Scientific Reasoning}

\begin{itemize}
\item \emph{Propositions} and \emph{Patterns of Proof}
\item Carroll, Lewis. ``What the tortoise said to Achilles.'' \emph{Mind} 104.416 (1995): 691-693.
\end{itemize}




\noindent \textsc{Oct 8: Sequences, Series and Limits}

\begin{itemize}
\item Pemberton \& Rau Ch. 5, and skim Ch. 31
\end{itemize}



\noindent \textsc{Oct 10:  Differentiation}

\begin{itemize}
\item Pemberton \& Rau Ch. 6, 7, appendix of Chapter 10

\end{itemize}



\noindent \textsc{Oct 15: Optimization}

\begin{itemize}
\item Pemberton \& Rau Ch. 8, 10

\end{itemize}


\noindent \textsc{Oct 17: Integration}

\begin{itemize}
\item Pemberton \& Rau Ch. 19, 20

\end{itemize}


\noindent \textsc{Oct 22:  Linear Algebra}

\begin{itemize}

\item Pemberton \& Rau Ch. 11, 12, 13.1


\end{itemize}


\noindent \textsc{Oct 24:  Multivariate Calculus}

\begin{itemize}

\item Pemberton \& Rau Ch. 13, 14, 15, 16

\end{itemize}
\noindent \textsc{Oct 29:  \underline{In-Class Midterm Exam}}



\newpage

\noindent \textsc{Oct 30: Laws and Properties of Probability}

\begin{itemize}

\item Bertsekas-Tsitsiklis Ch. 1.1-1.2 


\end{itemize}

\noindent \textsc{Nov 5: Conditional Probability}

\begin{itemize}

\item Bertsekas-Tsitsiklis Ch. 1.3-1.7

\end{itemize}

\noindent \textsc{Nov 7: Discrete Random variables }

\begin{itemize}

 \item Bertsekas-Tsitsiklis Ch. 2  
 


\end{itemize}


\noindent \textsc{Nov 12: General (Continuous) Random variables}

\begin{itemize}
\item Bertsekas-Tsitsiklis Ch. 3 

\end{itemize}


\noindent \textsc{Nov 14:  Summarizing Random variables}

\begin{itemize}
\item Bertsekas-Tsitsiklis Ch. 3.4-3.7 

\end{itemize}


\noindent \textsc{Nov 19:Covariance and Correlation}

\begin{itemize}
\item Bertsekas-Tsitsiklis Ch. 4.1-4.3

\end{itemize}

\noindent \textsc{Nov 21: Transformations of Random Variables }

\begin{itemize}

 \item Bertsekas-Tsitsiklis Ch. 4.4-4.4.6
 
 
\end{itemize}

\noindent \textsc{Nov 26:  Inequalities and limits of sequences of Random Variables}

\begin{itemize}

 \item  Bertsekas-Tsitsiklis Ch. 5


\end{itemize}

\noindent \textsc{Nov 28: Hypothesis testing}

\begin{itemize}

 \item Bertsekas-Tsitsiklis Ch. 9
  \item Degroot and Schervish (optional) Chapter 8

\end{itemize}

\noindent \textsc{Dec 3: Causal Inference}

\begin{itemize}

\item Angrist Pischke \emph{Mostly Harmless Econometrics} Sections 1, 2, 3.1 - 3.2 (pp 27 - 67), 4.1 - 4.2 (pp 113- 146), 4.6 (pp 188- 215)

\end{itemize}

\noindent \textsc{Dec 5: Review}



\end{document}














